\documentclass[11pt,a4paper]{article}
\usepackage[margin=1in]{geometry}
\usepackage{amsmath,amssymb,amsthm}
\usepackage{algorithm}
\usepackage{algpseudocode}
\usepackage{booktabs}
\usepackage{hyperref}

\theoremstyle{definition}
\newtheorem{definition}{\textbf{Definition}}[section]
\newtheorem{theorem}{\textbf{Theorem}}[section]
\newtheorem{lemma}[theorem]{\textbf{Lemma}}
\newtheorem{remark}{\textbf{Remark}}[section]

\title{\textbf{Sinusoidal Positional Encoding}}
\author{\textbf{Team Scaibu} \\ \small Email: \texttt{jaydeep.9664920749@gmail.com} \\ \small \textbf{Architecture and Core Engine Team}}
\date{\textbf{October 2026}}

\begin{document}
\maketitle

\section{Definitions and Cognitive Mental Models}

\subsection{Formal Mathematical Definition}
\textbf{Definition (Sinusoidal Positional Encoding Operator):}
Let $N \in \mathbb{N}_{\ge 1}$ denote maximum sequence length, $d_{\text{model}} \in 2\mathbb{N}_{\ge 1}$ denote an even model embedding dimension, and $B \in \mathbb{R}_{> 1}$ denote the geometric frequency base (typically $B = 10000.0$).
The \textbf{Sinusoidal Positional Encoding} operator:
{\boldmath
\begin{equation}
\operatorname{PE}: \mathbb{N}_{\ge 1} \times 2\mathbb{N}_{\ge 1} \times \mathbb{R}_{> 1} \longrightarrow \mathbb{R}^{N \times d_{\text{model}}}
\end{equation}
}
is the deterministic matrix-valued mapping whose matrix elements for token position $p \in \{0, \dots, N-1\}$ and frequency channel index $i \in \{0, \dots, d_{\text{model}}/2 - 1\}$ are defined by:
{\boldmath
\begin{equation}
\operatorname{PE}_{p, 2i} = \sin\left(\frac{p}{B^{2i / d_{\text{model}}}}\right), \qquad \operatorname{PE}_{p, 2i+1} = \cos\left(\frac{p}{B^{2i / d_{\text{model}}}}\right)
\end{equation}
}

\subsection{Expanded Non-Technical Definition (Intuition for Anyone)}
\textbf{Intuitive Conceptual Definition:}
Unlike Recurrent Neural Networks (RNNs) that process sentences word-by-word in sequential order, Transformers process all words in a sentence simultaneously in parallel. While this parallel processing enables massive GPU acceleration, it introduces a critical mathematical blindspot: **permutation invariance**. To a raw Transformer attention layer, \emph{``dog bites man''} and \emph{``man bites dog''} look 100\% identical because attention treats sequences as unordered sets of vectors.

Sinusoidal Positional Encoding solves this by stamping a unique, non-learned geometric fingerprint onto every word position:
\begin{enumerate}
    \item \textbf{A Multi-Speed Clockwork System}: The embedding dimensions act as a massive collection of synchronized clocks ticking at different speeds. The lowest dimensions oscillate with rapid, short wavelengths (ticking every 1 to 2 words, like a fast second hand), while the highest dimensions oscillate with vast wavelengths (spanning over 10,000 words, like an hour hand or calendar).
    \item \textbf{Continuous Geometric Uniqueness}: Every single position in the document receives an unmistakably unique combination of sine and cosine values across its $d_{\text{model}}$ coordinates.
    \item \textbf{Relative Distance as Linear Rotation}: Because of trigonometric angle-addition identities ($\cos(\alpha - \beta) = \cos \alpha \cos \beta + \sin \alpha \sin \beta$), a neural network can determine the relative distance $k$ between two words at positions $p$ and $p+k$ through a simple, fixed linear transformation, without needing to relearn position representations for every new distance.
\end{enumerate}

\subsection{Expanded Mental Model for Visualization (Understandable by a Child)}
\textbf{The Magic Cosmic Clock Tower with Hundreds of Pendulums}:
Imagine you are inside a towering magical clock tower with $d_{\text{model}}/2$ pendulums swinging from the wooden ceiling.

\begin{itemize}
    \item \textbf{The Pendulums of Many Speeds (Frequencies $\omega_i$)}:
    Every pendulum is a different length:
    \begin{itemize}
        \item Pendulum 0 is tiny like a hummingbird wing and swings back and forth super fast: \emph{tick-tick-tick-tick!} Every single word makes it complete a full swing.
        \item Pendulum 1 is slightly longer and swings half as fast.
        \item The giant Pendulum at the back of the tower is so huge that it takes 10,000 words just to swing from one side of the room to the other!
    \end{itemize}
    
    \item \textbf{The Snapshot Code (Sine and Cosine Coordinates)}:
    Every time a word steps onto the page, you take a photograph of all the pendulums at that exact second:
    \begin{itemize}
        \item The horizontal position of pendulum $i$ is its \textbf{Sine}.
        \item The vertical height of pendulum $i$ is its \textbf{Cosine}.
    \end{itemize}
    
    \item \textbf{The Unbreakable Fingerprint}:
    When word 0 enters, all pendulums are at their starting points.
    When word 1 enters, the tiny humming pendulums have swung all the way over, while the giant pendulum has barely moved a millimeter.
    No two words in the entire library ever have the same pattern of pendulum angles!
    
    \item \textbf{Measuring Distance by Angle}:
    If you want to know how far apart Word A and Word B are, you don't have to count words one by one—you just look at how much the pendulums rotated between them!
\end{itemize}

\section{Core Mathematical Equations: Formal Definitions, Meaning, and Importance}

\subsection{\textbf{Equation 1: Geometric Angular Frequency Progression}}
{\boldmath
\begin{equation}
\omega_i = \frac{1}{B^{2i / d_{\text{model}}}} = \exp\left( -\frac{2i}{d_{\text{model}}} \log B \right), \qquad i \in \left\{0, 1, \dots, \frac{d_{\text{model}}}{2} - 1\right\}
\end{equation}
}
\begin{itemize}
    \item \textbf{Formal Mathematical Definition}:
    The geometric progression mapping that spaces angular frequencies logarithmically across the interval $[B^{-1}, 1.0]$.
    \item \textbf{What It Means}:
    Wavelengths $\lambda_i = 2\pi / \omega_i$ form a geometric spectrum ranging from $2\pi \approx 6.28$ tokens up to $2\pi \times 10000 \approx 62,831$ tokens.
    \item \textbf{Why It Is Important}:
    Logarithmic frequency spacing guarantees that the model can resolve fine-grained local syntax (adjacent tokens) and broad document-level structure (thousands of tokens apart) simultaneously across different channels.
\end{itemize}

\subsection{\textbf{Equation 2: Even-Indexed Harmonic Coordinate}}
{\boldmath
\begin{equation}
\operatorname{PE}_{p, 2i} = \sin(p \cdot \omega_i) = \sin\left(\frac{p}{B^{2i / d_{\text{model}}}}\right)
\end{equation}
}
\begin{itemize}
    \item \textbf{Formal Mathematical Definition}:
    The odd-symmetric harmonic basis component evaluated at discrete position $p$ for subspace channel $i$.
    \item \textbf{What It Means}:
    The sine component starts at $\sin(0) = 0$ and oscillates bounded strictly within $[-1.0, 1.0]$.
    \item \textbf{Why It Is Important}:
    Sine coordinates provide odd parity ($\sin(-x) = -\sin(x)$), enabling the model to learn directional orientation (distinguishing whether token $k$ appears before or after token $p$).
\end{itemize}

\subsection{\textbf{Equation 3: Odd-Indexed Orthogonal Quadrature Coordinate}}
{\boldmath
\begin{equation}
\operatorname{PE}_{p, 2i+1} = \cos(p \cdot \omega_i) = \cos\left(\frac{p}{B^{2i / d_{\text{model}}}}\right)
\end{equation}
}
\begin{itemize}
    \item \textbf{Formal Mathematical Definition}:
    The even-symmetric harmonic quadrature basis component phase-shifted by $\pi/2$ radians relative to the even coordinate.
    \item \textbf{What It Means}:
    The cosine component starts at $\cos(0) = 1.0$ and pairs with the sine component to form a 2D circular phasor: $(\sin(p\omega), \cos(p\omega))$.
    \item \textbf{Why It Is Important}:
    Pairing sine with orthogonal cosine forms a complete complex exponential basis $e^{j p \omega_i}$, ensuring that the 2D norm of each frequency channel is strictly constant ($\sin^2 + \cos^2 \equiv 1$).
\end{itemize}

\subsection{\textbf{Equation 4: Relative Shift Linear Rotation Matrix}}
{\boldmath
\begin{equation}
\begin{bmatrix} \operatorname{PE}_{p+k, 2i} \\ \operatorname{PE}_{p+k, 2i+1} \end{bmatrix} = \begin{bmatrix} \cos(k \omega_i) & \sin(k \omega_i) \\ -\sin(k \omega_i) & \cos(k \omega_i) \end{bmatrix} \begin{bmatrix} \operatorname{PE}_{p, 2i} \\ \operatorname{PE}_{p, 2i+1} \end{bmatrix} = \mathbf{R}(k \omega_i) \begin{bmatrix} \operatorname{PE}_{p, 2i} \\ \operatorname{PE}_{p, 2i+1} \end{bmatrix}
\end{equation}
}
\begin{itemize}
    \item \textbf{Formal Mathematical Definition}:
    The special orthogonal rotation operator $\mathbf{R}(k \omega_i) \in \mathrm{SO}(2)$ that shifts 2D positional coordinates by offset $k$ via matrix multiplication.
    \item \textbf{What It Means}:
    The encoding for position $p+k$ is a linear transformation of the encoding at position $p$, where the transformation matrix depends only on the relative distance $k$, not on the absolute position $p$.
    \item \textbf{Why It Is Important}:
    This proves mathematically that a Transformer can learn to attend to relative positions ($k$ tokens away) using standard linear projection weights, enabling seamless length extrapolation.
\end{itemize}

\subsection{\textbf{Equation 5: Inner Product Kernel Distance Decay}}
{\boldmath
\begin{equation}
\langle \mathbf{pe}_p, \mathbf{pe}_{p+k} \rangle = \sum_{i=0}^{d_{\text{model}}/2 - 1} \left( \sin(p\omega_i)\sin((p+k)\omega_i) + \cos(p\omega_i)\cos((p+k)\omega_i) \right) = \sum_{i=0}^{d_{\text{model}}/2 - 1} \cos(k \omega_i)
\end{equation}
}
\begin{itemize}
    \item \textbf{Formal Mathematical Definition}:
    The canonical Euclidean dot product between the positional vectors of two tokens separated by relative displacement $k$.
    \item \textbf{What It Means}:
    The dot product depends strictly on relative offset $k$, achieving its maximum value $d_{\text{model}}/2$ at $k=0$ and decaying smoothly as relative distance $|k|$ grows.
    \item \textbf{Why It Is Important}:
    This inherent geometric decay naturally induces a soft locality bias in dot-product attention, allowing models to easily focus on nearby tokens.
\end{itemize}

\section{Rigorous Mathematical Analysis Checklist}

\subsection{[ ] Notation: Symbol and Operator Specification}
\begin{itemize}
    \item $N \in \mathbb{N}_{\ge 1}$: Sequence length (positions $p \in \{0, \dots, N-1\}$).
    \item $d_{\text{model}} \in 2\mathbb{N}_{\ge 1}$: Model embedding dimension (must be even).
    \item $B \in \mathbb{R}_{> 1}$: Frequency base (default $10000.0$).
    \item $\omega_i \in \mathbb{R}_{> 0}$: Angular frequency for channel index $i \in \{0, \dots, d_{\text{model}}/2 - 1\}$.
    \item $\operatorname{PE} \in \mathbb{R}^{N \times d_{\text{model}}}$: Full positional encoding matrix.
    \item $\mathbf{pe}_p \in \mathbb{R}^{d_{\text{model}}}$: Positional encoding vector for token position $p$.
    \item $\mathbf{R}(\theta) \in \mathrm{SO}(2)$: 2D rotation matrix of angle $\theta$.
\end{itemize}

\subsection{[ ] Definitions: Epistemic Nature of Variables}
\begin{table}[h]
\centering
\caption{\textbf{Epistemic Classification of Mathematical Quantities}}
\begin{tabular}{lll}
\toprule
\textbf{Variable} & \textbf{Epistemic Classification} & \textbf{Mathematical Nature} \\
\midrule
$N, d_{\text{model}}, B$ & \textbf{Defined} (Architectural parameters) & Natural / positive real numbers \\
$\omega_i$ & \textbf{Calculated} (Fixed frequency schedule) & Strictly decreasing geometric sequence \\
$\operatorname{PE}_{p, 2i}, \operatorname{PE}_{p, 2i+1}$ & \textbf{Calculated} (Deterministic trigonometric values) & Bounded reals in $[-1.0, 1.0]$ \\
$\operatorname{PE}$ & \textbf{Calculated} (Static lookup matrix) & Real matrix in $\mathbb{R}^{N \times d_{\text{model}}}$ \\
\bottomrule
\end{tabular}
\end{table}

\subsection{[ ] Structure: Composite Algebraic Operator Graph}
{\boldmath
\begin{equation}
(N, d_{\text{model}}, B) \xrightarrow{\text{Freq Schedule}} \{\omega_i\} \xrightarrow{\text{Outer Product } \mathbf{p} \otimes \boldsymbol{\omega}} \mathbf{\Theta} \xrightarrow{(\sin, \cos)} \operatorname{PE}
\end{equation}
}

\subsection{[ ] Units and Dimensions: Dimensional Consistency}
{\boldmath
\begin{align}
p &\in \{0, \dots, N-1\} \quad \text{\textbf{(Dimensionless discrete position)}} \\
\omega_i &\in \mathbb{R}_{> 0} \quad \text{\textbf{(Radians per token step)}} \\
p \cdot \omega_i &\in \mathbb{R}_{\ge 0} \quad \text{\textbf{(Phase angle in radians)}} \\
\operatorname{PE} &\in \mathbb{R}^{N \times d_{\text{model}}} \quad \text{\textbf{(Bounded coordinate tensor)}}
\end{align}
}

\subsection{[ ] Behavior: Limiting Regimes and Parameter Scaling}
\begin{itemize}
    \item \textbf{Position Zero ($p = 0$)}: $\operatorname{PE}_{0, 2i} = \sin(0) = 0$ and $\operatorname{PE}_{0, 2i+1} = \cos(0) = 1$ for all channels $i$.
    \item \textbf{High Frequency Limit ($i = 0$)}: $\omega_0 = 1.0$. The fastest channel has period $\lambda_0 = 2\pi \approx 6.28$ tokens.
    \item \textbf{Low Frequency Limit ($i = d_{\text{model}}/2 - 1$)}: $\omega_{\text{min}} \approx 1/B = 10^{-4}$. The slowest channel has period $\lambda \approx 62,831$ tokens.
\end{itemize}

\subsection{[ ] Conditions and Failure Cases}
\begin{itemize}
    \item \textbf{Odd Dimension Exception}: If $d_{\text{model}} \pmod 2 \ne 0$, sine and cosine channels cannot pair equally. \emph{Precondition}: $d_{\text{model}} \ge 2$ and even.
    \item \textbf{Zero or Negative Length}: Sequence length must satisfy $N \ge 1$.
\end{itemize}

\subsection{[ ] Verification and Invariant Properties}
\begin{itemize}
    \item \textbf{Constant Vector L2 Norm}: For every position $p \in \{0, \dots, N-1\}$, the Euclidean norm of the positional vector is strictly invariant:
    {\boldmath
    \begin{equation}
    \|\mathbf{pe}_p\|_2^2 = \sum_{i=0}^{d_{\text{model}}/2 - 1} \left( \sin^2(p\omega_i) + \cos^2(p\omega_i) \right) = \frac{d_{\text{model}}}{2} \implies \|\mathbf{pe}_p\|_2 = \sqrt{\frac{d_{\text{model}}}{2}}
    \end{equation}
    }
\end{itemize}

\subsection{[ ] Transfer: Harmonic Analysis & Fourier Feature Mapping}
Sinusoidal positional encoding is mathematically identical to Random Fourier Feature (RFF) coordinate mapping used in Neural Radiance Fields (NeRF) and kernel ridge regression, projecting 1D coordinates onto a multidimensional torus $\mathbb{T}^{d/2} \subset \mathbb{R}^d$.

\section{Theorems, Proofs, and Theoretical Guarantees}

\begin{theorem}[\textbf{Exact Linear Relative Shift Invariance}]
For any fixed integer displacement $k \in \mathbb{Z}$, there exists an orthogonal linear transformation $\mathbf{M}_k \in \mathrm{SO}(d_{\text{model}})$ such that for all positions $p \ge 0$:
{\boldmath
\begin{equation}
\mathbf{pe}_{p+k} = \mathbf{M}_k \mathbf{pe}_p
\end{equation}
}
\end{theorem}

\begin{proof}
Consider the 2D subspace corresponding to frequency channel $i$ with coordinates $(x_p, y_p) = (\sin(p\omega_i), \cos(p\omega_i))$.
Applying the standard trigonometric angle sum and difference identities:
{\boldmath
\begin{align}
\sin((p+k)\omega_i) &= \sin(p\omega_i)\cos(k\omega_i) + \cos(p\omega_i)\sin(k\omega_i) \\
\cos((p+k)\omega_i) &= \cos(p\omega_i)\cos(k\omega_i) - \sin(p\omega_i)\sin(k\omega_i)
\end{align}
}
Expressing this linear combination in matrix form:
{\boldmath
\begin{equation}
\begin{bmatrix} \sin((p+k)\omega_i) \\ \cos((p+k)\omega_i) \end{bmatrix} = \begin{bmatrix} \cos(k\omega_i) & \sin(k\omega_i) \\ -\sin(k\omega_i) & \cos(k\omega_i) \end{bmatrix} \begin{bmatrix} \sin(p\omega_i) \\ \cos(p\omega_i) \end{bmatrix} = \mathbf{R}_i(k) \begin{bmatrix} \sin(p\omega_i) \\ \cos(p\omega_i) \end{bmatrix}
\end{equation}
}
Notice that $\det(\mathbf{R}_i(k)) = \cos^2(k\omega_i) + \sin^2(k\omega_i) = 1.0$, so $\mathbf{R}_i(k) \in \mathrm{SO}(2)$.
Assembling the block-diagonal matrix $\mathbf{M}_k = \operatorname{diag}(\mathbf{R}_0(k), \mathbf{R}_1(k), \dots, \mathbf{R}_{d/2-1}(k))$, we have $\mathbf{pe}_{p+k} = \mathbf{M}_k \mathbf{pe}_p$, which is completely independent of absolute position $p$.
\end{proof}

\begin{theorem}[\textbf{Strict L2 Norm Conservation Invariant}]
For any dimension $d_{\text{model}} \in 2\mathbb{N}$ and any position $p \in \mathbb{N}$, the Euclidean norm satisfies:
{\boldmath
\begin{equation}
\|\mathbf{pe}_p\|_2 = \sqrt{\frac{d_{\text{model}}}{2}}
\end{equation}
}
\end{theorem}

\begin{proof}
Evaluating the squared Euclidean norm directly:
{\boldmath
\begin{equation}
\|\mathbf{pe}_p\|_2^2 = \sum_{j=0}^{d_{\text{model}}-1} (\operatorname{PE}_{p, j})^2 = \sum_{i=0}^{d_{\text{model}}/2 - 1} \left( \sin^2(p\omega_i) + \cos^2(p\omega_i) \right)
\end{equation}
}
By the fundamental Pythagorean trigonometric identity, $\sin^2(\theta) + \cos^2(\theta) \equiv 1$ for all $\theta \in \mathbb{R}$.
Therefore, the sum simplifies to:
{\boldmath
\begin{equation}
\|\mathbf{pe}_p\|_2^2 = \sum_{i=0}^{d_{\text{model}}/2 - 1} 1 = \frac{d_{\text{model}}}{2} \implies \|\mathbf{pe}_p\|_2 = \sqrt{\frac{d_{\text{model}}}{2}}
\end{equation}
}
Because this norm is identical for all positions $p$, positional encodings contribute equal geometric energy across all tokens without biasing early or late positions.
\end{proof}

\section{Critical Questions, Meta-Questions, and Deep Understanding}

\subsection{\textbf{Critical Question 1: Addition vs Concatenation of Positional Vectors}}
\textbf{Question}: Why does the standard Transformer add positional vectors to word embeddings ($\mathbf{x} + \mathbf{pe}$) rather than concatenating them ($[\mathbf{x} \,\|\, \mathbf{pe}]$)?
\begin{itemize}
    \item \textbf{Meta-Question}: \emph{How does linear projection in subsequent layers preserve subspace orthogonality even under additive vector superposition?}
    \item \textbf{Deep Answer}: Concatenating $\mathbf{x} \in \mathbb{R}^d$ and $\mathbf{pe} \in \mathbb{R}^d$ would double the embedding dimension to $2d$, increasing the number of parameters and FLOPs in all subsequent attention projection matrices by a factor of 4. Adding them preserves dimension $d$. In high-dimensional spaces ($d \ge 512$), random or structured vectors occupy near-orthogonal subspaces. The subsequent query and key weight matrices $\mathbf{W}^Q, \mathbf{W}^K$ learn to project word semantics and positional coordinates onto distinct, non-interfering attention subspaces.
\end{itemize}

\subsection{\textbf{Critical Question 2: Why Base 10000.0?}}
\textbf{Question}: Why is the base frequency constant chosen as $B = 10000.0$ in the original architecture?
\begin{itemize}
    \item \textbf{Meta-Question}: \emph{What governs the maximum non-repeating positional horizon under geometric frequency scaling?}
    \item \textbf{Deep Answer}: The maximum period in the encoding is $\lambda_{\text{max}} = 2\pi B \approx 62,831$ tokens. Setting $B = 10000$ guarantees that the slowest-moving channels do not complete even a single full oscillation over typical context windows (e.g. 512 to 4,096 tokens). If $B$ were too small (e.g. 100), the slowest channel would wrap around every 628 tokens, causing distant positions to produce duplicate encoding vectors and confusing the model.
\end{itemize}

\subsection{\textbf{Critical Question 3: Sinusoidal vs Learned Absolute Positional Embeddings}}
\textbf{Question}: Why do models like BERT use learned positional embeddings while models like Transformer-Vaswani use fixed sinusoidal encodings?
\begin{itemize}
    \item \textbf{Meta-Question}: \emph{What is the inductive bias trade-off between parametric expressivity and sequence length extrapolation?}
    \item \textbf{Deep Answer}: Learned positional embeddings ($\mathbf{W}_{\text{pos}} \in \mathbb{R}^{N_{\text{max}} \times d}$) treat each position as an unconstrained independent parameter, providing maximum flexibility to fit the training distribution. However, learned embeddings cannot extrapolate to positions beyond $N_{\text{max}}$ (position $N_{\text{max}}+1$ has no trained weights). Sinusoidal encodings have zero learnable parameters and can theoretically be evaluated at arbitrary sequence lengths $p \to \infty$ without retraining.
\end{itemize}

\subsection{\textbf{Critical Question 4: Transition from Sinusoidal to Rotary Embeddings (RoPE)}}
\textbf{Question}: Why have modern LLMs (e.g. LLaMA, Mistral) replaced additive sinusoidal encodings with Rotary Positional Embeddings (RoPE)?
\begin{itemize}
    \item \textbf{Meta-Question}: \emph{Why is multiplicative complex rotation in attention logits superior to additive input superposition?}
    \item \textbf{Deep Answer}: Additive encodings ($\mathbf{x} + \mathbf{pe}$) bleed positional noise directly into the semantic feature space, which must pass through LayerNorm and nonlinear MLP blocks, degrading relative distance precision. RoPE applies the exact same sinusoidal frequencies $\omega_i$, but applies them multiplicatively as 2D rotations directly to Query and Key vectors inside each attention head. In RoPE, the dot product $\mathbf{q}_m^\top \mathbf{k}_n$ becomes an exact function of relative distance $(m-n)$, eliminating positional noise from the value pathway completely.
\end{itemize}

\section{Algorithmic Specification}

The computational pipeline implemented in \texttt{NnAlgoSinusoidalPositionalEncoding} is formalized in Algorithm~\ref{alg:sinusoidal_pe}.

\begin{algorithm}[H]
\caption{Sinusoidal Positional Encoding Generation}
\label{alg:sinusoidal_pe}
\begin{algorithmic}[1]
\Require Sequence length $N \in \mathbb{N}_{\ge 1}$, embedding dimension $d_{\text{model}} \in 2\mathbb{N}_{\ge 1}$, base $B \in \mathbb{R}_{> 1}$
\Ensure Positional encoding matrix $\operatorname{PE} \in \mathbb{R}^{N \times d_{\text{model}}}$
\State $\textbf{Assert } N \ge 1 \textbf{ and } d_{\text{model}} \ge 2 \textbf{ and } d_{\text{model}} \pmod 2 = 0$
\State $d_{\text{half}} \gets d_{\text{model}} / 2$
\State $\operatorname{PE} \gets \text{zeros of shape } (N, d_{\text{model}})$
\State $\boldsymbol{\omega} \gets \text{array of size } d_{\text{half}}$
\For{$i \gets 0 \textbf{ to } d_{\text{half}} - 1$}
    \State $\omega_i \gets \exp\left( -\frac{2 \cdot i}{d_{\text{model}}} \cdot \log B \right)$
\EndFor
\For{$p \gets 0 \textbf{ to } N - 1$}
    \For{$i \gets 0 \textbf{ to } d_{\text{half}} - 1$}
        \State $\theta \gets p \cdot \omega_i$
        \State $\operatorname{PE}[p, 2 \cdot i] \gets \sin(\theta)$ \Comment{Even channels receive sine}
        \State $\operatorname{PE}[p, 2 \cdot i + 1] \gets \cos(\theta)$ \Comment{Odd channels receive cosine}
    \EndFor
\EndFor
\State \Return $\{\text{encoding\_matrix}: \operatorname{PE}, \text{seq\_len}: N, \text{d\_model}: d_{\text{model}}\}$
\end{algorithmic}
\end{algorithm}

\section{Complexity Analysis}
\begin{itemize}
    \item \textbf{Time Complexity}:
    \begin{itemize}
        \item Angular frequency precomputation: $\Theta(d_{\text{model}}/2)$ operations.
        \item Harmonic matrix population: $N \times (d_{\text{model}}/2)$ sine evaluations and $N \times (d_{\text{model}}/2)$ cosine evaluations.
        \item Total Time Complexity: $\Theta(N \cdot d_{\text{model}})$ operations.
    \end{itemize}
    \item \textbf{Space Complexity}:
    \begin{itemize}
        \item Output matrix buffer $\operatorname{PE}$: $\Theta(N \cdot d_{\text{model}})$ floats.
        \item Total Auxiliary Space: $\Theta(N \cdot d_{\text{model}})$.
    \end{itemize}
\end{itemize}

\section{Repository Reference}
All mathematical formulations, algorithmic implementations, contracts, and test suites are maintained at:
\begin{center}
\url{https://github.com/Chief-Strategist-J/llm-obs-infra}
\end{center}

\end{document}
